% Migrated from Linguistic_transparency on 2026-06-14.
% Source: paper/sections/02-morphosyntactic-transparency.tex, §2.5 "Transparent free relatives".

A second morphosyntactic-transparency case fits the schema. Transparent free relatives (TFRs) are a sub-type of fused-relative construction \autocite[ch.~12, §6.3]{huddleston-pullum-2002} in which the matrix accesses properties of an embedded predicate (the \term{nucleus}) through \mention{what}: \mention{what they call ecstatic}, \mention{what we regard as essential}, \mention{what seems to be a tourist}. \mention{What} is a fused-relative pronoun, morphologically and syntactically a singular indefinite; the matrix accesses features of the nucleus rather than of \mention{what}. The construction has been examined extensively in the formal-syntactic literature, with three main competing analyses: a parenthetical-with-backward-deletion treatment \autocite{wilder-1999-transparent-free-relatives,schelfhout-coppen-oostdijk-2003-transparent-free-relatives}, a unified analysis on which standard and transparent free relatives share a configurational structure \autocite{grosu-2003-unified-theory-tfrs,grosu-2016-semantics-syntax-morphology-tfrs}, and a graft analysis in which a single constituent is shared between matrix and relative clauses \autocite{vanriemsdijk-2006-free-relatives}. The present treatment follows Kim's \autocite*[ch.~6.8]{kim-2026-form-function} construction-grammar analysis and draws on the corpus evidence summarized in Kim and Reynolds \autocite*{kim-reynolds-2026-tfr-two-kinds}.

Two $R$-relations dissociate empirically, and they have different verb-class profiles.

\ea\label{ex:tfr-r-relations}
\ea \mention{What appear to be trousers are on the chair.}
\ex \mention{What they call essential workers are exhausted.}
\ex \mention{He was what they call dimwitted.}
\z
\z

\term{Agreement transparency} matches matrix-verb agreement to the number profile of the nucleus and is licensed across the full inventory of TFR-licensing verbs, including the \mention{seem}/\mention{appear} class. The first example in (\ref{ex:tfr-r-relations}) has plural agreement on plural nucleus \mention{trousers} despite singular \mention{what}; the form is corpus-attested (see \autocite{kim-reynolds-2026-tfr-two-kinds}). Attributional verbs license the same, as the second example shows. \term{Syntactic-category transparency} is the narrower relation: the matrix's syntactic-category requirement is mediated through the syntactic category of the nucleus, so a non-NP nucleus licenses the TFR in a non-NP slot. The third example has an AP nucleus (\mention{dimwitted}) and the TFR in predicative-AP position. This narrower relation is licensed only by the attributional sub-group of TFR-licensing verbs (\mention{call, consider, describe~as, regard~as, take~to~be}); the \mention{seem}/\mention{appear} class shows agreement transparency only.\footnote{A reader scanning CGEL for fused-relative examples with \mention{seem} or \mention{appear} will find NP-nucleus cases, since CGEL ch.~12 §6.3 treats fused relatives generally rather than the TFR sub-type; those examples are agreement-transparency cases under the present analysis, consistent with the verb-class restriction holding only for syntactic-category transparency.}

A working corpus check supports the asymmetry. AP nuclei are productive with attributional verbs (estimated hundreds of hits in COCA for \mention{call} and \mention{consider} alone, with sampled verification; the working-note data include \mention{what I consider beautiful}, \mention{what I consider important}, and \mention{what he takes to be obvious}) but absent with \mention{seem} and \mention{appear} in a 178-hit exhaustive search; \mention{take~to~be} patterns with the attributional class despite carrying an overt copula, indicating the constraint is on argument structure (the predicative complement's subject identified with the matrix object) rather than the presence of \mention{be} \autocite{kim-reynolds-2026-tfr-two-kinds}. Within the attributional class, PP nuclei are restricted to characterizing PPs (\mention{what she'd call in love with Sam}), not locatives, suggesting a further sub-type stratification on the nucleus's predicative profile.

The projectibility analysis runs in two task settings, parallel to §\ref{sec:2-4}. For \term{known} TFR-licensing verbs, the verb's argument-structure type (subject identified with matrix subject vs. matrix object) is sufficient input: object-orientation is the diagnostic for the attributional daughter and predicts syntactic-category transparency in addition to the agreement transparency licensed by the parent. For \term{candidate} verbs whose TFR-licensing status is uncertain (\mention{prove to be}, \mention{turn out to be}, \mention{happen to be} are open cases), the same argument-structure diagnostic plus the AP-nucleus test isolates daughter membership without verb-by-verb stipulation. Two off-list projections are available. \term{Diachronic}: the attributional class should show stable membership over time rather than free recruitment from the broader characterizing-verb inventory, with \mention{take~to~be}'s occupying a small grammaticalization pathway as the kind of marginal addition the cluster predicts. \term{Cross-linguistic}: languages with parallel matrix-licensing constructions should exhibit comparable parent/daughter splits if the constraint is on argument-structure configuration rather than on language-particular lexical content. Neither is tested in the present paper. Their availability as predictions, beyond the input criteria, supports treating the daughter as a candidate projectibility profile rather than as a distributional grouping.

The verb-class restriction itself admits two readings, both of which predict the asymmetry; distinguishing them is a research question the present paper leaves open. On a \term{metalinguistic} reading, the attributional construction is quotative: \mention{call} and \mention{consider} attribute a label rather than directly predicating, and the attribution context relaxes the matrix's selectional pressure on the nucleus's syntactic category in a way unavailable to ordinary predication. On a \term{mention-mediated} reading, the nucleus is treated as a use-mention hybrid; the construction marks the nucleus as a citation form whose syntactic-category requirement is suspended. The two differ on what licenses the relaxation: speaker-attribution versus citation-marking. They make the same first-pass predictions; cancellability and projection tests on the resulting CIs would be diagnostic.

TFRs are a second morphosyntactic projectibility profile, with the cluster of agreement, syntactic-category, and verb-class facts running off constructional sub-type membership rather than off the lexical-class membership that stabilizes §\ref{sec:2-1}. The stabilizing structure is constructional inheritance with a verb-class restriction at the daughter level. The constructional analysis is taken up in §\ref{sec:3-5}; full empirical and analytical development is in joint work with Jong-Bok Kim, in preparation.

% Source: paper/sections/03-constructional-transparency.tex, §3.5 "Transparent free relatives revisited".

The TFR agreement/syntactic-category dissociation is constructional transparency in a tight sense: a constructional sub-type (the attributional daughter) has access conditions that the parent lacks, and the access condition tracks an inheritance relation rather than a per-item lexical property. The TFR-parent licenses agreement transparency across the full verb inventory; the TFR-attributional daughter licenses syntactic-category transparency in addition. Kim's \autocite*[ch.~6.8]{kim-2026-form-function} treatment is the architectural anchor; the corpus check in Kim and Reynolds \autocite*{kim-reynolds-2026-tfr-two-kinds} supplies the parent/daughter split's empirical content. The two competing readings of the daughter's licence (§\ref{sec:2-5}: metalinguistic-attribution vs. mention-mediated) are both constructional in §\ref{sec:3}'s sense; the work §\ref{sec:3-5} does is independent of which reading wins.

The TFR case gives §\ref{sec:3-3}'s network hypothesis a tractable instantiation, though not a confirmation. The daughter's predictions (verb-class restriction, PP-restriction to characterizing PPs, the \mention{take~to~be} placement that defeats a copula-as-blocker analysis) are stabilized by the construction's position in an inheritance network, rather than by any verb's lexical entry or by the matrix's selectional requirements. $R$-relativity is also visible: the same mediator is transparent for agreement (both daughters) and only partially transparent for syntactic category (daughter-only). The CGEL-vs-CxG choice flagged in §\ref{sec:3-4} reappears: an inheritance-hierarchy framework like SBCG (Kim's choice) makes the parent/daughter relation the central analytical move; a Goldbergian variant absorbs the same relation into a network of family-resemblance constructions.
